
\documentclass{article}
\usepackage{colortbl}
\usepackage{makecell}
\usepackage{multirow}
\usepackage{supertabular}

\begin{document}

\newcounter{utterance}

\twocolumn

{ \footnotesize  \setcounter{utterance}{1}
\setlength{\tabcolsep}{0pt}
\begin{supertabular}{c@{$\;$}|p{.15\linewidth}@{}p{.15\linewidth}p{.15\linewidth}p{.15\linewidth}p{.15\linewidth}p{.15\linewidth}}

    \# & $\;$A & \multicolumn{4}{c}{Game Master} & $\;\:$B\\
    \hline 

    \theutterance \stepcounter{utterance}  

    & & \multicolumn{4}{p{0.6\linewidth}}{\cellcolor[rgb]{0.9,0.9,0.9}{%
	\makecell[{{p{\linewidth}}}]{% 
	  \tt {\tiny [A$\langle$GM]}  
	 State giocando a un gioco di negoziazione cooperativa in cui dovete accordarvi su come dividere una serie di oggetti tra voi e un altro giocatore.\\ \tt \\ \tt Regole:\\ \tt (a) A Lei e all'altro giocatore viene dato un insieme di oggetti. Ognuno di voi riceve anche una funzione di valore segreta, che rappresenta il valore di ciascun tipo di oggetto.\\ \tt (b) Si scambiano messaggi con l'altro giocatore per concordare chi si aggiudica gli oggetti. Potete inviare un massimo di 5 messaggi ciascuno, oppure terminare in anticipo facendo una proposta segreta in qualsiasi momento.\\ \tt (c) A ciascuno di voi viene chiesto di inviare una proposta segreta indicando gli oggetti che desiderate, formattata tra parentesi quadre come segue: "[Proposta: <numero> <nome oggetto>, <numero> <nome oggetto>, <...>]".\\ \tt (d) Se le vostre proposte sono complementari, cioè ci sono abbastanza oggetti per soddisfare entrambe le proposte, a ciascun giocatore viene assegnato un punteggio basato sulla somma dei valori degli oggetti ricevuti. In caso contrario, entrambi ricevono zero punti.\\ \tt (e) L'obiettivo è massimizzare la somma del proprio punteggio e di quello dell'altro giocatore.\\ \tt \\ \tt Cominciamo.\\ \tt \\ \tt L'insieme degli oggetti disponibili è:\\ \tt \\ \tt 2 ganci, 2 piatti, 1 ciotola\\ \tt \\ \tt La funzione valore segreta è:\\ \tt \\ \tt gancio: 3, piatto: 1, ciotola: 2\\ \tt \\ \tt IMPORTANTE: i vostri messaggi, a meno che non si tratti di una proposta segreta, vengono trasmessi direttamente all'altro giocatore, quindi non includete alcuna risposta alle regole o testo di annuncio del vostro messaggio. Per fare una proposta segreta, utilizzate il formato indicato. Non utilizzare delle parentesi quadre quando si comunica all'altro giocatore, altrimenti verrà interpretata come una proposta segreta.\\ \tt  
	  } 
	   } 
	   } 
	 & \\ 
 

    \theutterance \stepcounter{utterance}  

    & \multicolumn{4}{p{0.6\linewidth}}{\cellcolor[rgb]{0.8,1,0.9}{%
	\makecell[{{p{\linewidth}}}]{% 
	  \tt {\tiny [A$\rangle$GM]}  
	 Ciao! Ecco i miei valori: ogni gancio vale 3, ogni piatto 1, la ciotola 2. Quali sono i tuoi valori? Possiamo puntare a una divisione che massimizzi la somma totale. Ad esempio, se per te alcuni oggetti valgono poco, li prendo io e tu prendi quelli che per te valgono di più. Qual è il tuo ordine di preferenza? 
	  } 
	   } 
	   } 
	 & & \\ 
 

    \theutterance \stepcounter{utterance}  

    & & & \multicolumn{2}{p{0.3\linewidth}}{\cellcolor[rgb]{0.95,0.95,0.95}{%
	\makecell[{{p{\linewidth}}}]{% 
	  \tt {\tiny [GM$|$GM]}  
	 continue 
	  } 
	   } 
	   } 
	 & & \\ 
 

    \theutterance \stepcounter{utterance}  

    & & & \multicolumn{2}{p{0.3\linewidth}}{\cellcolor[rgb]{0.95,0.95,0.95}{%
	\makecell[{{p{\linewidth}}}]{% 
	  \tt {\tiny [GM$|$GM]}  
	 Ciao! Ecco i miei valori: ogni gancio vale 3, ogni piatto 1, la ciotola 2. Quali sono i tuoi valori? Possiamo puntare a una divisione che massimizzi la somma totale. Ad esempio, se per te alcuni oggetti valgono poco, li prendo io e tu prendi quelli che per te valgono di più. Qual è il tuo ordine di preferenza? 
	  } 
	   } 
	   } 
	 & & \\ 
 

    \theutterance \stepcounter{utterance}  

    & & \multicolumn{4}{p{0.6\linewidth}}{\cellcolor[rgb]{0.9,0.9,0.9}{%
	\makecell[{{p{\linewidth}}}]{% 
	  \tt {\tiny [GM$\rangle$B]}  
	 State giocando a un gioco di negoziazione cooperativa in cui dovete accordarvi su come dividere una serie di oggetti tra voi e un altro giocatore.\\ \tt \\ \tt Regole:\\ \tt (a) A Lei e all'altro giocatore viene dato un insieme di oggetti. Ognuno di voi riceve anche una funzione di valore segreta, che rappresenta il valore di ciascun tipo di oggetto.\\ \tt (b) Si scambiano messaggi con l'altro giocatore per concordare chi si aggiudica gli oggetti. Potete inviare un massimo di 5 messaggi ciascuno, oppure terminare in anticipo facendo una proposta segreta in qualsiasi momento.\\ \tt (c) A ciascuno di voi viene chiesto di inviare una proposta segreta indicando gli oggetti che desiderate, formattata tra parentesi quadre come segue: "[Proposta: <numero> <nome oggetto>, <numero> <nome oggetto>, <...>]".\\ \tt (d) Se le vostre proposte sono complementari, cioè ci sono abbastanza oggetti per soddisfare entrambe le proposte, a ciascun giocatore viene assegnato un punteggio basato sulla somma dei valori degli oggetti ricevuti. In caso contrario, entrambi ricevono zero punti.\\ \tt (e) L'obiettivo è massimizzare la somma del proprio punteggio e di quello dell'altro giocatore.\\ \tt \\ \tt Cominciamo.\\ \tt \\ \tt L'insieme degli oggetti disponibili è:\\ \tt \\ \tt 2 ganci, 2 piatti, 1 ciotola\\ \tt \\ \tt La funzione valore segreta è:\\ \tt \\ \tt gancio: 1, piatto: 0, ciotola: 8\\ \tt \\ \tt IMPORTANTE: i vostri messaggi, a meno che non si tratti di una proposta segreta, vengono trasmessi direttamente all'altro giocatore, quindi non includete alcuna risposta alle regole o testo di annuncio del vostro messaggio. Per fare una proposta segreta, utilizzate il formato indicato. Non utilizzare delle parentesi quadre quando si comunica all'altro giocatore, altrimenti verrà interpretata come una proposta segreta.\\ \tt \\ \tt \\ \tt Ciao! Ecco i miei valori: ogni gancio vale 3, ogni piatto 1, la ciotola 2. Quali sono i tuoi valori? Possiamo puntare a una divisione che massimizzi la somma totale. Ad esempio, se per te alcuni oggetti valgono poco, li prendo io e tu prendi quelli che per te valgono di più. Qual è il tuo ordine di preferenza? 
	  } 
	   } 
	   } 
	 & \\ 
 

    \theutterance \stepcounter{utterance}  

    & & & \multicolumn{4}{p{0.6\linewidth}}{\cellcolor[rgb]{1,0.85,0.72}{%
	\makecell[{{p{\linewidth}}}]{% 
	  \tt {\tiny [GM$\langle$B]}  
	 Per me: gancio 1, piatto 0, ciotola 8. Quindi la ciotola è di gran lunga la più preziosa per me; i piatti non valgono nulla; i ganci valgono poco.\\ \tt \\ \tt Dato che per te i ganci valgono 3 ciascuno e la ciotola solo 2, mentre per me la ciotola vale 8, la somma totale massima si ottiene se:\\ \tt - Io prendo la ciotola (valore totale: 8 per me, 2 per te se la prendessi tu).\\ \tt - Tu prendi i 2 ganci (valore 6 per te, 2 per me se li prendessi io).\\ \tt - I piatti non contano per me e valgono 1 ciascuno per te, quindi ha senso che li prenda tu.\\ \tt \\ \tt Proposta cooperativa: tu prendi 2 ganci e 2 piatti, io prendo 1 ciotola. Somma totale: 8 (mia) + 6+2 (tua) = 16, che sembra il massimo. Ti va bene? 
	  } 
	   } 
	   } 
	  \\ 
 

    \theutterance \stepcounter{utterance}  

    & & & \multicolumn{2}{p{0.3\linewidth}}{\cellcolor[rgb]{0.95,0.95,0.95}{%
	\makecell[{{p{\linewidth}}}]{% 
	  \tt {\tiny [GM$|$GM]}  
	 continue 
	  } 
	   } 
	   } 
	 & & \\ 
 

    \theutterance \stepcounter{utterance}  

    & & & \multicolumn{2}{p{0.3\linewidth}}{\cellcolor[rgb]{0.95,0.95,0.95}{%
	\makecell[{{p{\linewidth}}}]{% 
	  \tt {\tiny [GM$|$GM]}  
	 Per me: gancio 1, piatto 0, ciotola 8. Quindi la ciotola è di gran lunga la più preziosa per me; i piatti non valgono nulla; i ganci valgono poco.\\ \tt \\ \tt Dato che per te i ganci valgono 3 ciascuno e la ciotola solo 2, mentre per me la ciotola vale 8, la somma totale massima si ottiene se:\\ \tt - Io prendo la ciotola (valore totale: 8 per me, 2 per te se la prendessi tu).\\ \tt - Tu prendi i 2 ganci (valore 6 per te, 2 per me se li prendessi io).\\ \tt - I piatti non contano per me e valgono 1 ciascuno per te, quindi ha senso che li prenda tu.\\ \tt \\ \tt Proposta cooperativa: tu prendi 2 ganci e 2 piatti, io prendo 1 ciotola. Somma totale: 8 (mia) + 6+2 (tua) = 16, che sembra il massimo. Ti va bene? 
	  } 
	   } 
	   } 
	 & & \\ 
 

    \theutterance \stepcounter{utterance}  

    & & \multicolumn{4}{p{0.6\linewidth}}{\cellcolor[rgb]{0.9,0.9,0.9}{%
	\makecell[{{p{\linewidth}}}]{% 
	  \tt {\tiny [A$\langle$GM]}  
	 Per me: gancio 1, piatto 0, ciotola 8. Quindi la ciotola è di gran lunga la più preziosa per me; i piatti non valgono nulla; i ganci valgono poco.\\ \tt \\ \tt Dato che per te i ganci valgono 3 ciascuno e la ciotola solo 2, mentre per me la ciotola vale 8, la somma totale massima si ottiene se:\\ \tt - Io prendo la ciotola (valore totale: 8 per me, 2 per te se la prendessi tu).\\ \tt - Tu prendi i 2 ganci (valore 6 per te, 2 per me se li prendessi io).\\ \tt - I piatti non contano per me e valgono 1 ciascuno per te, quindi ha senso che li prenda tu.\\ \tt \\ \tt Proposta cooperativa: tu prendi 2 ganci e 2 piatti, io prendo 1 ciotola. Somma totale: 8 (mia) + 6+2 (tua) = 16, che sembra il massimo. Ti va bene? 
	  } 
	   } 
	   } 
	 & \\ 
 

    \theutterance \stepcounter{utterance}  

    & \multicolumn{4}{p{0.6\linewidth}}{\cellcolor[rgb]{0.8,1,0.9}{%
	\makecell[{{p{\linewidth}}}]{% 
	  \tt {\tiny [A$\rangle$GM]}  
	 Mi va benissimo: massimizza la somma totale.\\ \tt \\ \tt [Proposta: 2 ganci, 2 piatti] 
	  } 
	   } 
	   } 
	 & & \\ 
 

    \theutterance \stepcounter{utterance}  

    & & & \multicolumn{2}{p{0.3\linewidth}}{\cellcolor[rgb]{0.95,0.95,0.95}{%
	\makecell[{{p{\linewidth}}}]{% 
	  \tt {\tiny [GM$|$GM]}  
	 proposal 
	  } 
	   } 
	   } 
	 & & \\ 
 

    \theutterance \stepcounter{utterance}  

    & & & \multicolumn{2}{p{0.3\linewidth}}{\cellcolor[rgb]{0.95,0.95,0.95}{%
	\makecell[{{p{\linewidth}}}]{% 
	  \tt {\tiny [GM$|$GM]}  
	 [2, 2, 0] 
	  } 
	   } 
	   } 
	 & & \\ 
 

    \theutterance \stepcounter{utterance}  

    & & \multicolumn{4}{p{0.6\linewidth}}{\cellcolor[rgb]{0.9,0.9,0.9}{%
	\makecell[{{p{\linewidth}}}]{% 
	  \tt {\tiny [GM$\rangle$B]}  
	 L'altro giocatore ha presentato una proposta segreta. Ora tocca a Lei presentare una sua proposta.\\ \tt  
	  } 
	   } 
	   } 
	 & \\ 
 

    \theutterance \stepcounter{utterance}  

    & & & \multicolumn{4}{p{0.6\linewidth}}{\cellcolor[rgb]{1,0.85,0.72}{%
	\makecell[{{p{\linewidth}}}]{% 
	  \tt {\tiny [GM$\langle$B]}  
	 [Proposta: 1 ciotola] 
	  } 
	   } 
	   } 
	  \\ 
 

    \theutterance \stepcounter{utterance}  

    & & & \multicolumn{2}{p{0.3\linewidth}}{\cellcolor[rgb]{0.95,0.95,0.95}{%
	\makecell[{{p{\linewidth}}}]{% 
	  \tt {\tiny [GM$|$GM]}  
	 proposal 
	  } 
	   } 
	   } 
	 & & \\ 
 

    \theutterance \stepcounter{utterance}  

    & & & \multicolumn{2}{p{0.3\linewidth}}{\cellcolor[rgb]{0.95,0.95,0.95}{%
	\makecell[{{p{\linewidth}}}]{% 
	  \tt {\tiny [GM$|$GM]}  
	 [0, 0, 1] 
	  } 
	   } 
	   } 
	 & & \\ 
 

    \theutterance \stepcounter{utterance}  

    & & & \multicolumn{2}{p{0.3\linewidth}}{\cellcolor[rgb]{0.95,0.95,0.95}{%
	\makecell[{{p{\linewidth}}}]{% 
	  \tt {\tiny [GM$|$GM]}  
	 [[2, 2, 0], [0, 0, 1]] 
	  } 
	   } 
	   } 
	 & & \\ 
 

\end{supertabular}
}

\end{document}
